% -*- coding: utf-8 -*-
% 问题三两阶段求解流程图（服务日入列 + 事件驱动时刻 + 覆盖兜底 + 逐日阶段二）
% v2：增加按天解耦结构、拆分候选列池生成步骤、修正阶段一目标描述
\begin{figure}[htbp]
\centering
\begin{tikzpicture}[
    box/.style={draw, rounded corners=3pt, minimum width=3.4cm, minimum height=0.75cm,
        align=center, font=\footnotesize, fill=blue!8},
    boxg/.style={draw, rounded corners=3pt, minimum width=3.4cm, minimum height=0.75cm,
        align=center, font=\footnotesize, fill=green!10},
    boxo/.style={draw, rounded corners=3pt, minimum width=3.4cm, minimum height=0.75cm,
        align=center, font=\footnotesize, fill=orange!12},
    boxr/.style={draw, rounded corners=3pt, minimum width=3.4cm, minimum height=0.75cm,
        align=center, font=\footnotesize, fill=red!8},
    boxp/.style={draw, rounded corners=3pt, minimum width=3.4cm, minimum height=0.75cm,
        align=center, font=\footnotesize, fill=purple!10},
    boxteal/.style={draw, rounded corners=3pt, minimum width=3.4cm, minimum height=0.75cm,
        align=center, font=\footnotesize, fill=teal!10},
    arrow/.style={->, >=stealth, thick, draw=black},
    darrow/.style={->, >=stealth, thick, draw=black, dashed},
]

% ==================== 阶段一标题 ====================
\node[font=\small\bfseries, fill=green!15, draw, rounded corners, minimum width=14cm, minimum height=0.5cm] (s1) at (0, 6.8) {阶段一：固定排班（3840 人，不含临时任务）};

% ==================== 第一行：时间窗解析与候选序列 ====================
\node[boxg] (window) at (-5.4, 5.7) {时间窗解析\\可行日期集 $D_p$ + CPM 收缩};
\node[boxg] (seq)    at (-1.4, 5.7) {候选序列枚举\\直达 / 设施对 / 穿梭 OD 对};
\node[boxg] (event)  at ( 3.0, 5.7) {事件驱动时刻\\窗口最早时刻 $\cup$\{06:00\}，上限 12 个};
\node[boxp] (cover)  at ( 7.2, 5.7) {覆盖兜底\\未覆盖人员生成专属直达列};

% ==================== 第二行：服务日入列 + 装载 + 剪枝 ====================
\node[boxg] (svcday) at (-4.2, 4.0) {服务日入列\\$p$ 在 $D_p$ 每天均生成候选列};
\node[boxg] (load)   at ( 0.0, 4.0) {EDF 双序满员装载\\非临时优先 + 临时填残座};
\node[box]  (prune)  at ( 4.2, 4.0) {列池剪枝\\(40,10) $\to$ 约 7 万列};

% ==================== 第三行：按天解耦 + 逐日 MIP ====================
\node[boxteal] (decouple) at (-3.6, 2.3) {按天解耦\\（定理 1）架次不跨天\\机队状态逐日重置};
\node[box] (mip) at ( 0.6, 2.3) {阶段一渐进式 MILP\\RelGap $0.05\!\to\!0.035$\\warm-start 早停};
\node[boxg]   (assign)    at ( 4.8, 2.3) {逐日飞机指派\\贪心按起飞时刻分配\\并发约束保证 0 失败};

% ==================== 阶段一输出 ====================
\node[boxg] (tstar) at (0, 1.0) {汇总 $T^* = \sum_{d=1}^{7} T^*_d$（每日日预算 $T^*_d$）};

% ==================== 阶段二标题 ====================
\node[font=\small\bfseries, fill=orange!15, draw, rounded corners, minimum width=14cm, minimum height=0.5cm] (s2) at (0, -0.4) {阶段二：临时任务加入（160 人，$T \le T^*$）};

% ==================== 阶段二流程（逐日） ====================
\node[boxo, below=0.4cm of s2, xshift=-3.6cm] (day_cols) {逐日列裁剪\\当日可行临时需求};
\node[boxo, below=0.4cm of s2]                 (prize)    {逐日 prize-collecting MIP\\$\max \sum z_p$，日预算 $T^*_d$};
\node[boxo, below=0.4cm of s2, xshift= 3.6cm] (rerun)    {阶段一架次作保底列\\注入 MIP start 恒有可行解};

% ==================== 最终输出 ====================
\node[boxr, below=0.6cm of prize] (output) {输出 q3-routes / q3-assignments\\（3840 全覆盖，临时 93/160 按预算承运）\\七类硬约束验证零违规};

% ==================== 主箭头 ====================
% 第一行内部
\draw[arrow] (window) -- (seq);
\draw[arrow] (seq)    -- (event);
\draw[arrow] (event)  -- (cover);

% 第一行 → 第二行
\draw[arrow] (window.south)  -- ++(0,-0.4) -| (svcday.north);
\draw[arrow] (event.south)   -- ++(0,-0.20) -| (load.north);
\draw[arrow] (cover.south)   -- ++(0,-0.4) -| (prune.north);

% 第二行内部
\draw[arrow] (svcday) -- (load);
\draw[arrow] (load)   -- (prune);

% 第二行 → 第三行
\draw[arrow] (prune.south)    -- ++(0,-0.3) -| (mip.north);
\draw[arrow] (svcday.south)   -- ++(0,-0.3) -| (decouple.north);

% 第三行内部
\draw[arrow] (decouple) -- (mip);
\draw[arrow] (mip)      -- (assign);

% 第三行 → 阶段一输出
\draw[arrow] (assign.south) |- (tstar.east);
\draw[arrow] (decouple.south) |- (tstar.west);

% 阶段一输出 → 阶段二
\draw[arrow] (tstar) -- (s2);

% 阶段二内部
\draw[arrow] (s2)       -- (day_cols);
\draw[arrow] (s2)       -- (prize);
\draw[arrow] (s2)       -- (rerun);
\draw[arrow] (day_cols) -- (prize);
\draw[arrow] (rerun)    -- (prize);

% 阶段二 → 输出
\draw[arrow] (prize) -- (output);

% ==================== 覆盖兜底反馈 ====================
\draw[darrow] (cover.south) -- ++(0,-0.4) -| (window.south);
\node[font=\scriptsize, align=center, anchor=east] at (-5.4, 5.05) {缺失人员$\to$补列};

\end{tikzpicture}
\caption{问题三两阶段求解流程（服务日入列 + 事件驱动时刻 + 覆盖兜底 + 逐日 prize-collecting）}
\label{fig:q3_flow}
\end{figure}
